% Working ACC 2027 manuscript draft.
% The article class is used so this source compiles without a conference-specific
% class file. Before submission, transfer the content into the current ACC
% author-kit class available from PaperPlaza Support and repeat the PDF check.
\documentclass[10pt,twocolumn]{article}

\usepackage[letterpaper,margin=0.72in,columnsep=0.23in]{geometry}
\usepackage{amsmath,amssymb,bm}
\usepackage{booktabs,array,multirow}
\usepackage{graphicx}
\usepackage{xcolor}
\usepackage{enumitem}
\usepackage[hidelinks]{hyperref}
\usepackage{microtype}

\graphicspath{{../results/reports/}}
\setlength{\parindent}{1em}
\setlength{\parskip}{0pt}
\setlist{nosep,leftmargin=*}

\newcommand{\sat}{\operatorname{sat}}
\newcommand{\clip}{\operatorname{clip}}
\newcommand{\sgn}{\operatorname{sgn}}
\newcommand{\R}{\mathbb{R}}
\newcommand{\T}{^{\mathsf T}}

\title{\textbf{Online PES-Adapted Neuromorphic Membrane--Gate Residual Control\\
for Actuator-Constrained UAV Quaternion Attitude Tracking}}

\author{Gandhi A. Contreras Torres}

\date{Working draft --- August 27, 2026}

\begin{document}
\maketitle

\begin{abstract}
This paper presents online-adaptive neuromorphic residual augmentation based on
coupled membrane--gate (CMG) dynamics for actuator-constrained quaternion attitude
architecture for actuator-constrained quaternion attitude control. A classical
quaternion proportional--derivative controller supplies a bounded stabilizing
torque command. Positive/negative demand pairs drive continuous CMG gate
states, and a projected Prescribed Error Sensitivity (PES) rule adapts a linear
decoder online. The decoded term is applied as a bounded additive residual, so
the classical command is retained rather than attenuated or replaced. All
canonical multirotor commands pass through the modeled actuator interface and,
for rotor-actuated platforms, an explicit collective-thrust/body-moment mixer.
In a 45-s MuJoCo Hard45 attitude mission, the online residual reduces mean
quaternion error on a Skydio X2 rotor-actuator model from 0.2941 to 0.1380
(53.1\%) with zero rotor saturation. In an explicitly labeled Crazyflie
four-rotor authority design study with 1.5$\times$ nominal per-rotor thrust, it
reduces the error from 0.2854 to 0.1323 (53.6\%), also without saturation. The
stock Crazyflie XML moment limits saturate for 99.88\% of the mission and cannot
support the same claim. Fixed- and adaptive-conductance CMG variants are
numerically indistinguishable in the present amplified-gate regime, whereas
online decoder adaptation accounts for the large improvement. A matched QUBE
Servo 3 QLabs experiment further shows that the same controller code executes
at 500 Hz with bounded voltage and zero reported faults, while providing
nominal accuracy equivalent to LQR. These results motivate a transparent
CMG--decoder--actuator co-design methodology and identify the
remaining robustness, frozen-weight, and hardware-validation requirements.
\end{abstract}

\section{Introduction}

Learning-based control is attractive when unmodeled dynamics, disturbances, or
operating-condition changes make a fixed controller conservative. However, a
neural controller that directly replaces a classical law can obscure the
stability backbone, make actuator feasibility difficult to interpret, and
confound learning performance with changes in effective feedback gain. This is
particularly consequential for multirotor attitude control, where a requested
body torque is useful only if it can be realized by nonnegative, bounded rotor
commands.

This work studies a different arrangement: a classical quaternion controller
remains responsible for the nominal closed loop, while a low-dimensional
neuromorphic dynamical system generates a bounded correction. We call its
two-state element a \emph{coupled membrane--gate (CMG) dynamical unit}: membrane
state $v$ and continuous gate state $r$ interact through a state-dependent
threshold and optional input-dependent conductance. We use ``unit'' rather than
claiming a new LIF neuron class because the applied controller output is $r$, not
a thresholded spike train. The
term \emph{neuromorphic} is used here in a precise architectural sense: control
demand is encoded as paired currents, processed by stateful neuron-inspired
dynamics, and decoded through local online learning. The applied output is a
continuous gate state rather than a counted spike train; no claim of event-only
computation or neuromorphic-hardware energy efficiency is made.

The principal contributions are:
\begin{enumerate}
    \item a baseline-preserving neuromorphic CMG residual controller with
    explicit channel bounds and actuator allocation;
    \item a clean experimental separation between adaptive CMG conductance
    and online PES decoder adaptation;
    \item a rotor-actuator-faithful Skydio X2 result with more than 50\% nominal
    Hard45 error reduction and zero actuator saturation;
    \item a transparent Crazyflie authority study that separates an infeasible
    stock virtual-moment interface from a declared four-rotor propulsion design;
    and
    \item a cross-backend QLabs transfer experiment that documents safe bounded
    execution and a negative nominal-performance result rather than overstating
    generality.
\end{enumerate}

The current draft intentionally restricts its main performance claims to the
multirotor simulations. QCar vision, physical 3DOF Hover control, and frozen-
weight generalization are future experiments and are not presented as completed
evidence.

\section{Architecture and Signal Contract}

\subsection{Platform-independent controller contract}

Let $m$ denote the number of bounded baseline control channels exposed to the
augmentor. At sample $k$, the platform controller supplies
\begin{equation}
    \bm u_{0,k}\in\R^m, \qquad
    \bm \ell\in\R_{>0}^m,
\end{equation}
where $\ell_i$ is the symmetric magnitude limit for channel $i$. The
neuromorphic augmentor contains $2m$ CMG units---one positive and one
negative unit per control channel---and a decoder
\begin{equation}
    \bm W_k\in\R^{m\times 2m}.
\end{equation}
It returns the saturated command $\bm u_k$, applied residual
$\bm\delta_k=\bm u_k-\bm u_{0,k}$, continuous gate state $\bm r_k$, and
amplified activity $\bm a_k$.

Table~\ref{tab:channels} makes the dimensional mapping explicit. Plant degrees
of freedom (DOFs), controlled DOFs, residual channels, neurons, and physical
actuators are not interchangeable quantities.

\begin{table*}[t]
\centering
\caption{Platform, channel, and actuator mapping in the present code base.}
\label{tab:channels}
\small
\begin{tabular}{p{2.2cm}p{2.1cm}p{2.6cm}p{1.6cm}p{1.7cm}p{3.4cm}}
\toprule
\shortstack[l]{Platform/\\backend} & Plant coordinates & Controlled baseline output & Residual channels $m$ & CMG units $2m$ & Actuator interface and status \\
\midrule
Skydio X2 / MuJoCo & 6 rigid-body DOF; position plus quaternion & Body torque $(\tau_x,\tau_y,\tau_z)$; vertical collective is classical & 3 & 6 & Four nonnegative rotor-thrust actuators through a $4\times4$ mixer; canonical result \\
CF2 rotor design / MuJoCo & 6 rigid-body DOF & Body torque $(\tau_x,\tau_y,\tau_z)$ & 3 & 6 & Four declared rotor-thrust actuators; design-study result \\
Stock CF2 / MuJoCo & 6 rigid-body DOF & Body torque $(\tau_x,\tau_y,\tau_z)$ & 3 & 6 & One collective and three abstract moment actuators; arbitrary XML limits; authority boundary \\
QUBE / QLabs & 2 mechanical DOF $(\theta,\alpha)$; four-state feedback & Motor voltage $v$ during balance & 1 & 2 & PAL \texttt{QubeServo3}, \texttt{hardware=0}; canonical QLabs transfer \\
3DOF Hover surrogate / MuJoCo & 3 rotational DOF represented by a quaternion & Body moment $(\tau_x,\tau_y,\tau_z)$ & 3 & 6 & Applied body moment in integration surrogate; physical four-motor adapter not implemented \\
QCar surrogate / MuJoCo & Planar $(x,y,\psi)$ & Steering-like yaw command $\delta$ & 1 & 2 & Forward force plus yaw moment surrogate; integration only \\
QCar hardware / Jetson & Vehicle plus camera-derived lane state & Steering command & 1 proposed & 2 proposed & Existing HSV+PID/\allowbreak SMC/\allowbreak NMPC runner; canonical CMG/SNN connection not yet completed \\
\bottomrule
\end{tabular}
\end{table*}

\subsection{End-to-end workflow}

Figure~\ref{fig:workflow} shows one controller sample. The baseline path is
always evaluated first. Neuromorphic processing operates on the bounded
baseline demand, not directly on raw actuator commands. For the multirotors,
collective thrust is computed separately by the vertical controller and is not
decoded by the CMG bank.

\begin{figure*}[t]
\centering
\setlength{\tabcolsep}{2.2pt}
\renewcommand{\arraystretch}{1.35}
\scriptsize
\begin{tabular}{c@{\ $\rightarrow$\ }c@{\ $\rightarrow$\ }c@{\ $\rightarrow$\ }c@{\ $\rightarrow$\ }c@{\ $\rightarrow$\ }c}
\fbox{\shortstack{reference\\$q_d$ or $x_d$}}
& \fbox{\shortstack{state/error\\$e_k,\ \omega_k$}}
& \fbox{\shortstack{classical law\\$\bm u_{0,k}$}}
& \fbox{\shortstack{normalize\\pair encode}}
& \fbox{\shortstack{$2m$ CMG\\gate states}}
& \fbox{\shortstack{decoder\\$\bm s_k=\bm W_k\bm a_k$}}\\[2mm]
\multicolumn{6}{c}{%
\fbox{\shortstack{baseline $\bm u_{0,k}$ plus decoded $\bm s_k$}}
$\ \rightarrow\ $
\fbox{\shortstack{bounded additive\\residual}}
$\ \rightarrow\ $
\fbox{\shortstack{actuator limits\\and allocation}}
$\ \rightarrow\ $
\fbox{\shortstack{plant\\MuJoCo/QLabs/hardware}}
$\ \rightarrow\ $
\fbox{\shortstack{measured state\\at $k+1$}}}\\[2mm]
\multicolumn{6}{c}{%
\fbox{\shortstack{$\bm e_k,\bm\omega_k,\bm a_k,\bm s_k$}}
$\ \dashrightarrow\ $
\fbox{\shortstack{projected PES\\update}}
$\ \dashrightarrow\ $
\fbox{\shortstack{$\bm W_{k+1}$ used by\\decoder at $k+1$}}}
\end{tabular}
\par\vspace{1mm}
\footnotesize
Measured plant state closes the physical feedback loop from the second row to
the state/error block in the first row. Dashed arrows denote learning only.
\caption{One-sample signal workflow. Solid arrows are the applied control path;
dashed arrows are online learning. The decoder first produces $\bm s_k$ from
$\bm W_k$, then PES updates $\bm W_{k+1}$; the new weights affect the next
sample, not the command already being applied.}
\label{fig:workflow}
\end{figure*}

The implemented order is:
\begin{enumerate}
    \item read the simulator or platform state;
    \item generate the reference and form the shortest-path state error;
    \item compute and saturate the classical command $\bm u_{0,k}$;
    \item normalize $\bm u_{0,k}$ and construct $2m$ pair currents;
    \item update CMG gate and membrane states;
    \item amplify and clip the continuous gate activity;
    \item decode $\bm s_k$ with the current weights $\bm W_k$;
    \item if online PES is enabled, compute the target and update
    $\bm W_{k+1}$;
    \item form the current additive residual using the previously decoded
    $\bm s_k$;
    \item saturate, allocate, and apply the actuator command; and
    \item advance the plant and log errors, commands, residuals, activity,
    saturation, parameters, model hash, and run configuration.
\end{enumerate}

This ordering is important: the current command is causal in $\bm W_k$, while
the learning update changes the subsequent command.

\section{Coupled Membrane--Gate Residual Controller}

\subsection{Pair-coded control demand}

The bounded baseline is
\begin{equation}
    \bm u_{0,k}=\sat_{\bm\ell}(\bm u^{\rm raw}_{0,k}),
\end{equation}
and its elementwise normalized demand is
\begin{equation}
    \bar{\bm u}_k=\bm u_{0,k}\oslash\bm\ell\in[-1,1]^m.
\end{equation}
Positive and negative demand are separated without discarding sign:
\begin{align}
    \bm I_k^+ &= I_b\bm 1+I_g\max(\bar{\bm u}_k,\bm 0),\\
    \bm I_k^- &= I_b\bm 1+I_g\max(-\bar{\bm u}_k,\bm 0),\\
    \bm I_k &= [\bm I_k^{+\T},\bm I_k^{-\T}]\T\in\R^{2m}.
\end{align}
The current implementation uses $I_b=0.75$ and $I_g=1.25$. Pair coding gives
the decoder an explicit directional basis and initializes to an interpretable
signed mapping.

\subsection{Coupled membrane--gate dynamics}

For unit $j$, membrane state $v_j$ and gate state $r_j\in[0,1]$ evolve as
\begin{align}
    \dot r_j &= \frac{\sigma\!\left(
    \alpha_h[v_j-(v_{\rm th}-\gamma r_j)]\right)-r_j}{\tau_r},
    \label{eq:gate}\\
    \dot v_j &= -\frac{v_j-v_{\rm rest}}{\tau}
    +\frac{I_j}{\tau}-c_j(v_j-v_{\rm reset}),
    \label{eq:membrane}
\end{align}
where
\begin{equation}
    c_j=\left(g_0+\mathbf{1}_{\rm adapt}g_1|I_j|\right)r_j^2.
\end{equation}
The state-dependent threshold $v_{\rm th}-\gamma r_j$ produces positive
feedback in the gate state, while $c_j$ adds a gate-dependent conductance.
Fixed conductance sets $\mathbf{1}_{\rm adapt}=0$; adaptive conductance sets it
to one. This neuronal adaptation is distinct from decoder learning.

Table~\ref{tab:hyperparams} lists the code defaults. These are currently shared
by all canonical platforms unless a runner explicitly overrides them.

\begin{table}[t]
\centering
\caption{Canonical CMG and decoder defaults.}
\label{tab:hyperparams}
\small
\begin{tabular}{lll}
\toprule
Symbol/code field & Value & Meaning \\
\midrule
$\tau$ & 0.05 s & membrane time constant \\
$\tau_r$ & 0.005 s & gate time constant \\
$v_{\rm rest}$ & 0 & resting state \\
$v_{\rm reset}$ & 0 & conductance reversal/reset \\
$v_{\rm th}$ & 1 & nominal threshold \\
$\alpha_h$ & 50 & sigmoid steepness \\
$\gamma$ & 2 & gate-threshold coupling \\
$g_0$ & 120 & base conductance gain \\
$g_1$ & 30 & adaptive scaling \\
$I_b$ & 0.75 & pair-current bias \\
$I_g$ & 1.25 & pair-current demand gain \\
$G$ & $10^5$ & post-gate amplification \\
$W_{\max}$ & 6 & decoder projection bound \\
\bottomrule
\end{tabular}
\end{table}

The numerical update is not a conventional spike/reset loop. First, forward
Euler updates the gate:
\begin{equation}
    r_{j,k+1}=\clip_{[0,1]}\left[r_{j,k}
    +\frac{\Delta t}{\tau_r}(\hat r_{j,k}-r_{j,k})\right],
\end{equation}
with
\begin{equation}
    \hat r_{j,k}=\sigma\!\left(
    \alpha_h[v_{j,k}-(v_{\rm th}-\gamma r_{j,k})]\right).
\end{equation}
Then, using $r_{j,k+1}$, the membrane equation is integrated exactly over the
sample under a zero-order-hold current. Define
\begin{align}
    c_{j,k+1} &= (g_0+\mathbf{1}_{\rm adapt}g_1|I_{j,k}|)r_{j,k+1}^2,\\
    \lambda_{j,k} &= \tau^{-1}+c_{j,k+1},\\
    d_{j,k} &= v_{\rm rest}/\tau+I_{j,k}/\tau+c_{j,k+1}v_{\rm reset}.
\end{align}
The code applies
\begin{equation}
    v_{j,k+1}=v_{j,k}e^{-\lambda_{j,k}\Delta t}
    +\frac{d_{j,k}}{\lambda_{j,k}}
    \left(1-e^{-\lambda_{j,k}\Delta t}\right).
\end{equation}
The continuous activity delivered to the decoder is
\begin{equation}
    \bm a_k=\clip_{[0,1]}(G\bm r_{k+1}).
\end{equation}
Although the dynamics are neuron inspired, $\bm a_k$ is a bounded continuous
gate signal. Spike events are not counted in the current controller.

\subsection{Fixed decoder and bounded additive residual}

The decoder is initialized as
\begin{equation}
    \bm W_0=[\bm I_m\; -\bm I_m],
\end{equation}
so each output initially subtracts its negative-demand unit from its positive-
demand unit. The signed neural output is
\begin{equation}
    \bm s_k=\clip_{[-1,1]}(\bm W_k\bm a_k),
\end{equation}
and is dimensionalized by the channel limits:
\begin{equation}
    \bm u^{\rm neural}_k=\bm s_k\odot\bm\ell.
\end{equation}
The canonical additive architecture is
\begin{align}
    \bm u_k &= \sat_{\bm\ell}\left(
    \bm u_{0,k}+\beta\bm u^{\rm neural}_k\right),
    \label{eq:additive}\\
    \bm\delta_k &= \bm u_k-\bm u_{0,k}.
\end{align}
Equation~\eqref{eq:additive} preserves the baseline before final saturation.
This differs from the historical neural-blending ablation,
\begin{equation}
    \bm u_k=(1-\beta)\bm u_{0,k}
    +\beta\bm u^{\rm neural}_k,
\end{equation}
which changes the classical feedback gain even when the neural path is
uninformative. The 50\%+ historical Skydio reproduction uses the latter form
and is not merged with the canonical actuator-faithful claims.

\subsection{Online PES decoder adaptation}

For the multirotors, the normalized learning target is
\begin{equation}
    \bm y_k^*=\clip_{[-1,1]}
    (-k_e\bm e_{q,k}-k_\omega^*\bm\omega_k),
    \label{eq:target}
\end{equation}
where $\bm e_q$ is the shortest-path quaternion vector error. Define the
decoder error
\begin{equation}
    \bm\varepsilon_k=\bm y_k^*-\bm s_k.
\end{equation}
The implemented projected PES-style update follows the local activity/error
structure of Prescribed Error Sensitivity learning \cite{bekolay2013pes}:
\begin{equation}
    \bm W_{k+1}=\Pi_{[-6,6]}\left(
    \bm W_k+\eta\bm\varepsilon_k\bm a_k\T
    -\lambda_W\bm W_k\right).
    \label{eq:pes}
\end{equation}
The outer product changes each output weight in proportion to output error and
current presynaptic activity. Projection prevents unbounded decoder growth;
leak weakly regularizes all weights. In the Hard45 experiments,
$\eta=8\times10^{-4}$, $k_e=8$, and $k_\omega^*=0$. The target therefore
teaches the decoder an error-directed correction rather than a copy of the
baseline torque.

PES is ``online'' because Eq.~\eqref{eq:pes} is evaluated at every control
sample on the evaluation trajectory. It is not evidence of zero-shot
generalization. Furthermore, the current implementation computes $\bm s_k$
before updating $\bm W_{k+1}$; the update first affects $\bm u_{k+1}$.

The QUBE target follows the verified voltage sign convention:
\begin{equation}
    y_k^*=\clip_{[-1,1]}\left(
    k_e\frac{\alpha_k}{20^\circ}
    +k_d\frac{\dot\alpha_k}{6\ {\rm rad/s}}\right).
\end{equation}
The same two-unit decoder machinery is used, but the platform state, baseline
law, actuator units, and target definition are platform specific.

\section{Quaternion Baseline and Actuator Connection}

\subsection{Quaternion error and body-torque baseline}

Let $q=[q_0,\bm q_v\T]\T$ be the measured unit quaternion and $q_d$ the
reference. The code forms
\begin{equation}
    q_e=q_d\otimes q^{-1}
\end{equation}
and selects the shortest path through
\begin{equation}
    \bm e_q=\sgn(q_{e,0})\operatorname{vec}(q_e).
\end{equation}
This shortest-path construction is consistent with established quaternion
multirotor control formulations \cite{fresk2013quaternion}. With diagonal
inertia $\bm J$, natural frequencies
$\bm\omega_n=[4,4,3.5]\T$ rad/s, and damping ratio $\zeta=0.95$, the code uses
elementwise gains
\begin{align}
    \bm K_q &= 2\,\operatorname{diag}(\bm J)
    \odot\bm\omega_n^2,\\
    \bm K_\omega &=2\zeta\,\operatorname{diag}(\bm J)
    \odot\bm\omega_n.
\end{align}
The bounded classical body torque is
\begin{equation}
    \bm\tau_{0,k}=\sat_{\bm\tau_{\lim}}
    (-\bm K_q\odot\bm e_{q,k}
    -\bm K_\omega\odot\bm\omega_k).
\end{equation}
For the canonical multirotor controller, $m=3$,
$\bm u_0=\bm\tau_0$, and $\bm\ell=\bm\tau_{\lim}$.

\subsection{Vertical collective channel}

Attitude augmentation does not learn collective thrust. Let $m_q$ denote total
vehicle mass, $z$ altitude, $v_z$ vertical speed, and $R_{33}$ the body vertical
axis projection. The implementation uses
\begin{equation}
    T=\clip\left[
    \frac{m_q(g+6(1-z)-4v_z)}{\max(R_{33},0.35)},
    0,2.5m_qg\right].
\end{equation}
Thus the applied rotor allocation receives one classical collective channel
and three baseline-plus-neural body-moment channels.

\subsection{Four-rotor allocation}

For rotor-actuated X2 and CF2-design models,
\begin{equation}
    \begin{bmatrix}T&\tau_x&\tau_y&\tau_z\end{bmatrix}\T
    =\bm B\bm f,
\end{equation}
where $\bm f=[f_1,f_2,f_3,f_4]\T$ and
\begin{equation}
    \bm B=
    \begin{bmatrix}
    1&1&1&1\\
    y_1&y_2&y_3&y_4\\
    -x_1&-x_2&-x_3&-x_4\\
    \kappa_1&\kappa_2&\kappa_3&\kappa_4
    \end{bmatrix}.
\end{equation}
The controller computes
\begin{equation}
    \bm f=\clip_{[\bm f_{\min},\bm f_{\max}]}
    \left(\bm B^{-1}
    [T,\bm\tau\T]\T\right).
\end{equation}
Rotor saturation is logged whenever any $f_i$ reaches either control bound.
Consequently, a body-torque improvement cannot silently bypass actuator
authority in the canonical rotor-path results.

The stock CF2 XML is different. It exposes collective thrust and three abstract
body-moment actuators. Their limits are described by the source model as
arbitrary. They are useful for an authority-boundary experiment but not as
identified motor limits.

\section{Models, Backends, and Fidelity Boundaries}

\subsection{Skydio X2 rotor model}

The X2 path loads the vendored MuJoCo Menagerie scene and addresses four
actuators named \texttt{thrust1}--\texttt{thrust4}. Site positions and actuator
yaw gears are read from the model at runtime to construct $\bm B$. The
neuromorphic residual acts only through the torque request that is mixed into
these rotor commands. Plant disturbances, when enabled, are applied separately
through MuJoCo's external-force field and are never used to realize the control
command.

\subsection{Crazyflie stock model and rotor-authority design}

The stock Menagerie Crazyflie 2 model has mass 0.027 kg and diagonal inertia
\begin{equation}
    \operatorname{diag}(\bm J)=
    [2.3951,2.3951,3.2347]\times10^{-5}\ {\rm kg\,m^2}.
\end{equation}
Its XML exposes one body-thrust and three moment actuators with approximately
$10^{-5}$ N\,m moment authority per axis. The historical learned runner imposed
minimum limits $[0.03,0.03,0.02]$ N\,m, which are respectively
$[3000,3000,2000]$ times the XML moment authority and
$[104.38,104.38,61.83]$ times the project's inertia-derived limits. Those
historical torque floors are not used for the canonical CF2 claim.

The separate CF2 rotor-design model retains the Menagerie mass, inertia,
geometry, and sensors, but replaces abstract moments with four explicit thrust
actuators. The declared assumptions are a 92-mm motor diagonal, 32.5-mm
projected X-arm coordinates, base maximum thrust 0.0875 N per rotor, and
alternating yaw reaction coefficient 0.005 m. The paper configuration applies
a 1.5$\times$ authority scale:
\begin{equation}
    f_{\max}=0.13125\ {\rm N/rotor},\qquad
    T_{\max}=0.525\ {\rm N},
\end{equation}
which gives modeled thrust-to-weight ratio $T_{\max}/(mg)\approx1.98$. This is
a controller--actuator design study, not an identified Crazyflie 2.1
propulsion model. Stronger propellers alone are not assumed sufficient; motor,
ESC, battery, thermal, yaw-drag, and structural constraints require hardware
identification.

\subsection{QLabs QUBE virtual twin}

The canonical QUBE runner uses a Quanser-compatible Python 3.11 environment and
constructs
\begin{center}
\texttt{QubeServo3(id=0, hardware=0, frequency=500, pendulum=1)}.
\end{center}
PAL connects to the QLabs QUBE Servo 3 pendulum virtual card on port 18923. The
runner reads arm/pendulum position and rate, motor current, amplifier-fault, and
stall signals, and writes one motor voltage. Because PAL zeros encoders at
construction, the runner first commands 0 V, detects rest from 1.5 s of position
stability, records the hanging raw angle $\alpha_h$, and defines
\begin{equation}
    \alpha=\operatorname{wrap}(\alpha_{\rm raw}-\alpha_h-\pi).
\end{equation}
The tachometer is not used for rest detection because the twin exhibits an
approximately 0.061 rad/s bias at rest.

Energy swing-up is always classical. During captured balance, the state
$[\theta,\alpha,\dot\theta,\dot\alpha]\T$ is controlled by
\begin{equation}
    v_0=\sat_{5\,\mathrm V}
    ([-1,26.72,-0.8269,2.214]x),
\end{equation}
and the one-channel neuromorphic residual is capped at 0.35 V
($\beta=0.35/5=0.07$). A \texttt{finally} block commands 0 V, sets the LED red,
disables the virtual amplifier, and closes the card. The runner hard-codes
\texttt{hardware=0}; these are virtual-twin, not physical-hardware results.

\subsection{Integration surrogates and future hardware paths}

The MuJoCo pendulum, 3DOF Hover, and QCar XML files are integration surrogates.
They verify that the common controller API, channel dimensions, saturation, and
logging operate across different state spaces. They are not calibrated digital
twins. The 3DOF Hover repository currently contains vendor MATLAB/Simulink/QUARC
models but no reviewed Python HIL adapter. The QCar hardware repository contains
a camera/HSV lane estimator and PID, SMC, image-space NMPC, and earlier
earlier membrane--gate controllers under the historical ``HyperLIF'' label. The
canonical CMG code and temporal-
contrast SNN perception module are not yet connected to the Jetson runner.

\section{Experimental Protocol}

\subsection{Missions and controller conditions}

The principal result uses the 45-s Hard45 reference
\begin{align}
\phi_d(t)&=0.55\sin(0.75t)+0.20\sin(2.10t+0.40),\\
\theta_d(t)&=0.45\sin(0.95t+\pi/3)+0.18\sin(1.75t),\\
\psi_d(t)&=0.65\sin(0.55t+\pi)+0.25\sin(1.35t+0.20).
\end{align}
All angles are in radians. The integration step is $\Delta t=5\times10^{-4}$ s.
The compared conditions are:
\begin{enumerate}
    \item classical quaternion PD;
    \item fixed-conductance CMG with fixed decoder;
    \item adaptive-conductance CMG with fixed decoder;
    \item fixed-conductance CMG with online PES decoder; and
    \item adaptive-conductance CMG with online PES decoder.
\end{enumerate}
For the reported Hard45 residual runs, $\beta=0.5$. All conditions use the same
model, reference, integration step, baseline gains, torque policy, collective
controller, mixer, and actuator limits within a platform.

\subsection{Metrics}

The instantaneous quaternion-angle error is
\begin{equation}
    e_q^{\angle}(k)=2\operatorname{atan2}(\|\operatorname{vec}(q_{e,k})\|,
    |q_{e,0,k}|).
\end{equation}
The headline metric is its sample mean. We additionally log maximum error,
mean torque magnitude, residual norm, mean gate state, altitude variation, and
actuator saturation fraction. The project composite score is a secondary
diagnostic:
\begin{equation}
J_{\rm comp}=\bar e_q^{\angle}+0.02\,\overline{\|\bm\tau\|}
+0.20\,f_{\rm sat}+0.05\,\sigma_z,
\end{equation}
where $f_{\rm sat}$ is actuator-saturation fraction and $\sigma_z$ is altitude
standard deviation.

The QUBE protocol accumulates exactly 20 s of actual balance for every
controller, with a 90-s wall-time timeout, and compares the final 5 s. Three
trials are run per controller. Capture time is logged but not treated as a
matched metric because QLabs does not reset the arm state between processes.

\subsection{Robustness and generalization protocol}

The prepared five-seed robustness matrix adds a 10-degree initial attitude,
sinusoidal plant disturbance torque, attitude and rate measurement noise, and
seeds 1--5. A final submission must also include a frozen-weight protocol:
train or adapt on one trajectory, freeze $\bm W$, reset the plant, and evaluate
on held-out frequencies, amplitudes, and initial conditions. Until that
experiment is complete, the paper claims online adaptation on the evaluation
trajectory rather than zero-shot generalization.

\section{Results}

\subsection{Skydio X2 rotor-actuator result}

Table~\ref{tab:hard45} reports the validated nominal Hard45 result. Fixed and
adaptive conductance produce identical values, while online PES reduces mean
quaternion error by 53.07\% and the existing composite score by approximately
52.04\%. No compared X2 condition reaches a rotor bound.

\begin{table*}[t]
\centering
\caption{Validated 45-s Hard45 results. Improvement is relative to the
classical controller on the same platform.}
\label{tab:hard45}
\small
\resizebox{\textwidth}{!}{%
\begin{tabular}{llrrrr}
\toprule
Platform & Controller & Mean quaternion error & Improvement & Composite score & Saturation \\
\midrule
X2 rotor model & Classical & 0.294082 & 0.00\% & 0.297747 & 0.00\% \\
& CMG fixed decoder, fixed conductance & 0.263890 & 10.27\% & 0.267175 & 0.00\% \\
& CMG fixed decoder, adaptive conductance & 0.263890 & 10.27\% & 0.267175 & 0.00\% \\
& CMG--PES, fixed conductance & 0.138007 & \textbf{53.07\%} & 0.142785 & 0.00\% \\
& CMG--PES, adaptive conductance & 0.138007 & \textbf{53.07\%} & 0.142785 & 0.00\% \\
\midrule
Stock CF2 XML & Classical & 2.458662 & 0.00\% & 2.672362 & 99.88\% \\
& CMG fixed decoder, fixed conductance & 2.516808 & $-2.36$\% & 2.730492 & 99.89\% \\
& CMG fixed decoder, adaptive conductance & 2.517485 & $-2.39$\% & 2.731169 & 99.89\% \\
\midrule
CF2 rotor design, 1.5$\times$ & Classical & 0.285449 & 0.00\% & 0.287104 & 0.00\% \\
& CMG fixed decoder, fixed conductance & 0.266050 & 6.80\% & 0.267663 & 0.00\% \\
& CMG fixed decoder, adaptive conductance & 0.266050 & 6.80\% & 0.267663 & 0.00\% \\
& CMG--PES, fixed conductance & 0.132319 & \textbf{53.65\%} & 0.134444 & 0.00\% \\
& CMG--PES, adaptive conductance & 0.132319 & \textbf{53.65\%} & 0.134444 & 0.00\% \\
\bottomrule
\end{tabular}
}
\end{table*}

\IfFileExists{../results/reports/x2_hard45_qerr.png}{
\begin{figure}[t]
\centering
\includegraphics[width=\columnwidth]{x2_hard45_qerr.png}
\caption{Skydio X2 Hard45 quaternion-angle error. The online decoder adapts on
the plotted evaluation trajectory.}
\label{fig:x2}
\end{figure}}{}

\subsection{Crazyflie actuator-authority result}

The stock XML result is dominated by actuator saturation and cannot distinguish
controller quality on Hard45. In contrast, the explicit four-rotor design at
1.5$\times$ authority has zero saturation and closely reproduces the X2
online-PES improvement. This does not validate a specific hardware upgrade; it
states the modeled authority required for the controller to realize its torque.

\IfFileExists{../results/reports/cf2_authority_sweep.png}{
\begin{figure}[t]
\centering
\includegraphics[width=\columnwidth]{cf2_authority_sweep.png}
\caption{CF2 rotor-authority sweep. The actuator multiplier is a declared
design variable, not an identified stock-hardware coefficient.}
\label{fig:authority}
\end{figure}}{}

\subsection{Conductance versus decoder adaptation}

The fixed/adaptive-conductance pairs are numerically identical on both
rotor-actuated platforms. Mean online-run gate state is approximately 0.963;
post-gate amplification therefore places much of the activity near clipping.
In that regime, changing membrane conductance has little influence on the
decoded basis. Online PES, however, changes $\bm W$ and produces the large
improvement. The current evidence supports the statement ``decoder adaptation
matters under this parameterization,'' not ``adaptive CMG conductance
causes the improvement.'' A gate-gain/current-bias sensitivity study is needed
to characterize regimes where conductance adaptation remains observable.

\subsection{Neuron-family comparison: evidence boundary}

Archived exploratory runs provide a useful hypothesis but are not a controlled
neuron-model ablation. In those runs, direct LIF increased mean quaternion error
relative to PD by 12.2\% on CF2 and 19.6\% on X2, and an Izhikevich implementation
increased X2 error by 18.6\%; the historical CMG implementation reduced error by
1.2\% and 6.4\%, respectively. However, CMG decoded the smooth gate state $r$,
LIF decoded normalized voltage, Izhikevich decoded sparse events, and the legacy
controller used a 50/50 neural blend. No saved ALIF artifact is currently
available. These observations motivate the CMG hypothesis but do not establish
an inherent superiority of its dynamics.

A decisive comparison must hold fixed the six signed channels, input currents,
activity normalization, additive residual law, initial decoder, PES parameters,
actuator path, tuning budget, missions, and seeds. Required CMG ablations decode
$r$, $v$, and $[v,r]$, set $\gamma=0$, remove gate-dependent conductance, and
replace $r$ by a bandwidth-matched low-pass state. LIF, ALIF, and Izhikevich must
receive equally dimensional continuous and event-filtered readouts. Claims of an
intrinsic CMG advantage are deferred until this matched study is complete.

\subsection{Matched QUBE QLabs transfer}

Table~\ref{tab:qube} reports the corrected matched experiment. All 15 runs are
balanced in the final 5 s with zero reported faults and stalls. Online-PES
performance is practically equivalent to LQR. An earlier fixed-wall-time
comparison appeared to show approximately 14\% improvement, but unequal
swing-up duration gave controllers different amounts of balance and learning
time. Matching 20 s of actual balance removed the apparent advantage. The
invalidated comparison is retained in the project audit but is not a result.

\begin{table}[t]
\centering
\caption{QUBE Servo 3 QLabs matched nominal result ($n=3$).}
\label{tab:qube}
\small
\begin{tabular}{lrr}
\toprule
Controller & Final-5-s RMS (deg) & vs. LQR \\
\midrule
LQR & $1.0581\pm0.0012$ & reference \\
Fixed residual & $1.0590\pm0.0020$ & $-0.08$\% \\
Adaptive residual & $1.1170\pm0.1003$ & $-5.57$\% \\
Online PES fixed & $1.0566\pm0.0011$ & $+0.15$\% \\
Online PES adaptive & $1.0570\pm0.0015$ & $+0.10$\% \\
\bottomrule
\end{tabular}
\end{table}

QLabs encoder resolution is approximately $2\pi/2048=0.00307$ rad, or
0.176 degrees per count, and the twin exhibits equilibrium limit-cycle branches
near 1.06 and 1.23 degrees RMS. This reinforces why repeated, reset-controlled
protocols are necessary before assigning small differences to a controller.

\section{Discussion}

\subsection{What is genuinely neuromorphic?}

The architecture is neuromorphic because paired currents drive stateful
nonlinear units whose gate and membrane variables obey neuron-inspired dynamics,
and because a local activity/error outer product adapts the decoder online. It
is not currently event driven end to end: sensors are sampled periodically,
gate states are continuous, and MuJoCo/QLabs advance on conventional time
steps \cite{pedersen2024nir,nourse2023sns}. A future event-camera implementation could convert temporal visual
contrast into sparse ON/OFF events and feed the same paired-channel principle,
but such a result is outside the present evidence.

\subsection{Why online PES can improve the drones}

The fixed decoder begins as a directional signed mapping of the baseline
demand. PES then reshapes that mapping according to tracking error, allowing
cross-channel combinations of the six activities. Because $\bm W$ is
$3\times6$, each output torque may learn from every positive/negative axis
activity rather than only its own pair. This permits coupled residual structure
while the final torque and rotor commands remain bounded. The improvement may
include trajectory-specific compensation; frozen-weight hold-out evaluation is
required to distinguish learned reusable structure from online trajectory
adaptation.

\subsection{Actuator feasibility is part of the claim}

A neural body torque is not a physical command until allocation maps it to
feasible actuators. The X2 result is strong because the same rotor mixer and
bounds are used for every controller and saturation remains zero. The stock CF2
result demonstrates the converse: when the actuator model is saturated almost
continuously, comparing controllers is not meaningful. The enhanced CF2 model
therefore answers a co-design question---how much declared thrust authority is
needed---rather than claiming that current stock hardware already achieves the
result.

\subsection{Limitations and prohibited interpretations}

The following limitations are explicit:
\begin{itemize}
    \item the advantage of CMG over LIF, ALIF, and Izhikevich is currently a
    hypothesis because the archived implementations did not share a matched
    feature/decoder contract;
    \item online PES learns on the evaluated trajectory; zero-shot and frozen-
    weight generalization are not yet shown;
    \item the nominal X2/CF2 results are deterministic and require the prepared
    disturbance/noise/initial-condition matrix;
    \item the CF2 rotor coefficients and 1.5$\times$ propulsion scale are design
    assumptions, not identified hardware parameters;
    \item the QUBE experiment is a QLabs virtual twin, not physical hardware,
    and demonstrates safe equivalence rather than improved nominal accuracy;
    \item the 3DOF Hover Python HIL adapter is not implemented; and
    \item the QCar temporal-contrast SNN exists as an integration baseline but
    is not yet connected to physical steering.
\end{itemize}

\section{Reproducibility and Execution Workflow}

The MuJoCo simulations use the engine described in \cite{todorov2012mujoco}.
The clean repository is
\begin{center}
\url{C:/Research/papers/ACC2027_Gate_Residual_Control}.
\end{center}
Its Python 3.13 virtual environment contains MuJoCo 3.12 and the canonical
\texttt{gra\_control} package. The principal launchers are
\texttt{run\_x2\_paper\_matrix.ps1},
\texttt{run\_cf2\_paper\_matrix.ps1},
\texttt{run\_cf2\_authority\_sweep.ps1}, and
\texttt{run\_robustness\_matrix.ps1}. Each run writes a trace, summary, and
manifest containing model and configuration provenance.

The QUBE runner is
\begin{center}
\url{C:/Research/pendulum/runners/twin_acc2027_gate_residual.py}.
\end{center}
It executes under the Quanser-compatible Python 3.11 environment, imports the
same \texttt{gra\_control.gate} source, and records SHA-256 hashes of both the
runner and controller. The exact environment and commands are documented in
\texttt{docs/EXECUTION\_AND\_CODE\_MAP.md} and
\texttt{docs/PENDULUM\_QLABS\_INTEGRATION.md}.

\section{Conclusion}

An online-adaptive neuromorphic CMG residual can improve aggressive
quaternion attitude tracking while retaining a bounded classical control
backbone and explicit actuator allocation. On the X2 rotor model and an
explicitly labeled CF2 propulsion-authority design, online PES reduces nominal
Hard45 quaternion error by approximately 53\% without rotor saturation. The
large improvement is associated with decoder adaptation, whereas adaptive CMG
conductance is a null ablation in the present amplified-gate regime.
The stock CF2 result shows that actuator authority must be evaluated together
with the controller. A matched QLabs pendulum experiment demonstrates safe
cross-backend execution but no nominal accuracy benefit. Full robustness,
frozen-weight hold-out, identified propulsion parameters, and carefully staged
hardware experiments remain necessary before broader generalization claims.

\begin{thebibliography}{9}
\bibitem{fresk2013quaternion}
E. Fresk and G. Nikolakopoulos, ``Full quaternion based attitude control for a
quadrotor,'' in \emph{Proc. European Control Conf.}, 2013, pp. 3864--3869,
doi: 10.23919/ECC.2013.6669617.
\bibitem{bekolay2013pes}
T. Bekolay, C. Kolbeck, and C. Eliasmith, ``Simultaneous unsupervised and
supervised learning of cognitive functions in biologically plausible spiking
neural networks,'' in \emph{Proc. 35th Annual Conf. Cognitive Science Society},
2013, pp. 169--174.
\bibitem{pedersen2024nir}
J. E. Pedersen, S. Abreu, M. Jobst, \emph{et al.}, ``Neuromorphic
intermediate representation: A unified instruction set for interoperable
brain-inspired computing,'' \emph{Nature Communications}, vol. 15, art. 8122,
2024, doi: 10.1038/s41467-024-52259-9.
\bibitem{nourse2023sns}
W. R. P. Nourse, C. Jackson, N. S. Szczecinski, and R. D. Quinn,
``SNS-Toolbox: An open source tool for designing synthetic nervous systems and
interfacing them with cyber--physical systems,'' \emph{Biomimetics}, vol. 8,
no. 2, art. 247, 2023, doi: 10.3390/biomimetics8020247.
\bibitem{todorov2012mujoco}
E. Todorov, T. Erez, and Y. Tassa, ``MuJoCo: A physics engine for model-based
control,'' in \emph{Proc. IEEE/RSJ Int. Conf. Intelligent Robots and Systems},
2012, pp. 5026--5033, doi: 10.1109/IROS.2012.6386109.
\end{thebibliography}

\end{document}
